\section{Separating the conformal data of the two ends}
\label{sh:section}
\label{sec:shape}

The estimates of Section~\ref{sp:section} are expressed in the
uniformizing coordinates of the entire annulus.  To use them in a
transfer operator, their leading term must instead depend separately
on the inner and outer explored regions.  We prove the requisite
comparison uniformly over their shapes.  Only separation of the
two ends is needed; no boundary smoothness bound is imposed.

Let $K_i$ be a closed Jordan region containing $0$ in its interior,
and let $U_o$ be a bounded Jordan domain containing $K_i$.
The annulus under consideration is $U_o\setminus K_i$.
Fix constants $c_*,C_*>0$ and assume that
\begin{equation}\label{sh:shape-hypotheses}
 K_i\subset\{|z|\le a\},\quad
 \operatorname{diam}K_i\ge c_*a,\qquad
 \{|z|<b\}\subset U_o,\quad
 \min_{z\in\partial U_o}|z|\le C_*b.
\end{equation}
We take $b/a$ larger than a constant depending only on $c_*,C_*$.
Each boundary carries four distinct marked prime ends and alternating
color prescriptions.

Normalize the two ends separately.  Let
\[
 w_i:\widehat\C\setminus K_i\longrightarrow\{|w|>1\},
 \qquad w_i(z)=z/l_i+O(1)\quad(z\to\infty),
\]
with $l_i>0$, and let
\[
 w_o:U_o\longrightarrow\{|w|<1\},\qquad
 w_o(0)=0,\quad w_o'(0)=1/l_o>0.
\]
Thus $l_i$ is the logarithmic capacity of $K_i$ and $l_o$ is the
conformal radius of $U_o$ at $0$.  Write $y^0,x^0$ for the angles of
the inner and outer marks in these maps.  Capacity-diameter bounds
for continua and the Koebe estimates imply
\begin{equation}\label{sh:end-scales}
 l_i\asymp a,\qquad l_o\asymp b,
\end{equation}
with constants depending only on $c_*,C_*$.

\begin{lemma}[Uniform end-map comparison]\label{sh:comparison}
Let $(p,x,y)$ be the modulus and marked angles in a uniformization
of $U_o\setminus K_i$, with angular direction the same on both
boundary components.  Put $p'=\log(l_o/l_i)$.  The common rotational
normalization of the annulus map can be chosen so that
\begin{equation}\label{sh:comparison-estimate}
 |p-p'|\le C\frac ab,\qquad
 |x-x^0|+|y-y^0|
       \le C\bigl(1+\log(b/a)\bigr)\frac ab.
\end{equation}
Distances in the second inequality use compatible cyclic lifts.
The constant is uniform over all Jordan boundaries satisfying
\eqref{sh:shape-hypotheses}, including all placements of the marks.
\end{lemma}

\begin{proof}
Define the single-valued analytic corrections
\[
 g_i(z)=\log\frac{l_iw_i(z)}z,
 \qquad g_o(z)=\log\frac{l_ow_o(z)}z,
\]
with branches vanishing at infinity and at zero, respectively.
The quotient defining $g_o$ has no zeros in $U_o$ and the apparent
singularity at zero is removable; the exterior quotient has zero
winding and its logarithm is normalized at infinity.
Univalent distortion and \eqref{sh:end-scales} give
\begin{equation}\label{sh:correction-bounds}
 |g_o(z)|\le Ca/b\quad (|z|\le C_1a),\qquad
 |g_i(z)|\le Ca/b\quad (|z|\ge c_1b),
\end{equation}
for every fixed $C_1<\infty$ and $c_1>0$, once $a/b$ is small
enough.  The relevant sets in these inequalities are contained in
the respective map domains.  The same bounds hold in fixed round
collars in the single-end coordinates.  If $\psi_i=w_i^{-1}$,
the maximum principle applied to $\psi_i(w)/w$ gives
$|\psi_i(w)|\le a|w|$ for $|w|>1$.
For $\psi_o=w_o^{-1}$, Koebe's distortion bound gives
$|\psi_o(w)|\ge l_o|w|/(1+|w|)^2$.
Thus the inverse exterior map sends $1\le|w_i|\le2$ into
$|z|\le2a$, and the inverse interior map sends, for example,
$1/2\le|w_o|<1$ into $|z|\ge cb$.
These observations are the reason arbitrary roughness of the physical
boundaries does not change the constants.

Consider the multivalued analytic function
\[
 H(z)=\log(z/l_i)+g_i(z)+g_o(z).
\]
It has period $2\pi i$, and its real part is within $Ca/b$ of $0$
on the inner boundary and of $p'$ on the outer boundary.  Let
$A(z)$ be an annulus logarithm, normalized by
$0<\operatorname{Re}A<p$.  Its period is also $2\pi i$.
The difference $H-A$ is therefore single-valued.  In the cylinder
coordinate $A$, the circular average of the real part of any
single-valued analytic function is constant across the cylinder.
At the two boundary circles the averages of
$\operatorname{Re}(H-A)$ are, respectively, $O(a/b)$ and
$p'-p+O(a/b)$.  Equality of the averages proves
$|p-p'|\le Ca/b$.  The maximum principle now yields
\begin{equation}\label{sh:real-error}
 |\operatorname{Re}(A-H)|\le Ca/b
 \quad\hbox{throughout }U_o\setminus K_i.
\end{equation}

It remains to compare the angular parts; a bound on real parts
alone would not control them on arbitrary boundaries.  Use instead
the round collars in the single-end coordinates.  The relevant
single-valued functions satisfy
\[
 A-\log w_i=(A-H)+g_o,\qquad
 A-p-\log w_o=(A-H)+(p'-p)+g_i.
\]
Near the inner
end, $A-\log w_i$ has real part bounded by $Ca/b$, by
\eqref{sh:correction-bounds} and \eqref{sh:real-error}, and its
real part is exactly zero on $|w_i|=1$.  Schwarz reflection across
this circle, followed by interior gradient bounds on a smaller fixed
collar, bounds the variation of its imaginary part along the entire
end circle by $Ca/b$.  Near the outer end apply the same argument
to $A-p-\log w_o$.  Its real part is zero on $|w_o|=1$ and is
bounded by $Ca/b$ in a fixed inner collar.  This gives the analogous
angular bound, initially with possibly different additive constants
at the two ends.

To match these constants, choose $C_1$ large and $c_1$ small, both
fixed.  The round annulus $C_1a<|z|<c_1b$ is contained in the
physical annulus.  Along a radial segment across it, use
$O(1+\log(b/a))$ overlapping disks of fixed relative size.
The gradient estimate for the analytic function $A-H$, whose real
part obeys \eqref{sh:real-error}, bounds the variation of its
imaginary part on each disk by $Ca/b$.
At either end, distortion in the corresponding single-end chart
connects this radial segment to the fixed collar with a bounded
additional chart distance.  The same interior estimate and the
collar reflection estimate apply there.
Thus the two additive constants differ by at most
$C(1+\log(b/a))a/b$.  Choose the common rotation to fix the inner
constant.  The comparison of all marked angles, uniformly even
when marks are very close, follows.  This proves
\eqref{sh:comparison-estimate}.
\end{proof}

\begin{proposition}[Separated-end transmission]\label{sh:separated-ends}
There are constants $\sigma_1,\sigma_2>0$ and $C<\infty$,
depending only on the fixed constants in
\eqref{sh:shape-hypotheses}, such that the transmission probabilities
of every such marked annulus satisfy
\begin{equation}\label{sh:rank-one}
 \left|T^\pm-
 K\,l_i^\lambda J(y^0)^{1/3}\,
       l_o^{-\lambda}J(x^0)^{1/3}\right|
 \le C(a/b)^{\lambda+\sigma_1}.
\end{equation}
Here $K$ is the common constant in
Proposition~\ref{sp:transmission-asymptotics}.  If one rotates the
inner end, with its marks and prescribed colors, through angle $\phi$
about $0$ while keeping the outer end fixed, then
\begin{equation}\label{sh:rotation-average}
 \left|\frac1{2\pi}\int_0^{2\pi}
                 \bigl(T^+_\phi-T^-_\phi\bigr)\,d\phi\right|
       \le C(a/b)^{2+\sigma_2}.
\end{equation}
The same conclusion holds if the outer end is rotated instead.
Both estimates remain uniform for Jordan approximants to the
prime-end boundaries of Section~\ref{sec:annuli}, provided they
satisfy \eqref{sh:shape-hypotheses} with the same fixed constants.
\end{proposition}

\begin{proof}
Write $\rho=a/b$ and $q=\log(1/\rho)$.  By
\eqref{sh:end-scales}, $p'=q+O(1)$.  Set
$\eta=C(1+q)\rho$, the angular error in
Lemma~\ref{sh:comparison}.  The $1/3$-H\"older bound for $h$
and the first inequality of that lemma give
\[
 |e^{-\lambda p}h(x,y)-e^{-\lambda p'}h(x^0,y^0)|
       \le C\rho^\lambda(\rho+\eta^{1/3}).
\]
Combining this with \eqref{sp:rank-one-round} proves
\eqref{sh:rank-one}; for example any
$0<\sigma_1<\min\{\sigma,1/3\}$ is sufficient after changing
the constant.  Since
$e^{-\lambda p'}=(l_i/l_o)^\lambda$, the leading term has exactly
the displayed separated form.

For the signed estimate put $U=T^+-T^-=e^{-\lambda p}hW_-$.
Equations~\eqref{sp:odd-decay} and
\eqref{sp:odd-regularity}, together with the H\"older bound for
$h$, show that replacing the actual annular data by the two
single-end data costs at most
\begin{equation}\label{sh:signed-shape-error}
 C\rho^2\bigl(\rho^{1/2}+\eta+\eta^{1/3}\bigr).
\end{equation}
This bound is uniform over the rotation of either physical end.
In particular it is a power improvement on $\rho^2$.

The capacity $l_i$ is unchanged by rotation of the inner end, and
its normalized map is
$w_{i,\phi}(z)=e^{i\phi}w_i(e^{-i\phi}z)$.
Consequently its marked single-end angles are exactly
$y^0+\phi\mathbf1$.  The factor $h$ is unchanged by rotating
either tuple separately.  Thus the rotation average in the unperturbed
data is
\[
 e^{-\lambda p'}h(x^0,y^0)
                       (\Pi W_-)(p',x^0,y^0),
\]
whose absolute value is at most
$C(1+q)\rho^{\lambda+1}=C(1+q)\rho^{9/4}$ by
\eqref{sp:odd-decay}.  Add the uniform error
\eqref{sh:signed-shape-error}.
Any $0<\sigma_2<1/4$ is then admissible, after increasing the
constant and restricting to sufficiently small $\rho$.
Rotating the outer end gives the same projection, by simultaneous
rotation invariance.

Every estimate was obtained with constants depending only on the
separation and scale conditions, not on boundary smoothness or mark
separation.  It can therefore be applied directly to the continuum
function $T$ evaluated on the conformal data of the Jordan cycles
approximating an explored domain.  Under admissible convergence,
Proposition~\ref{an:limit} separately identifies the limit of the
discrete transmission probabilities with that continuum function; no
reflection across a physically nonsimple limiting bank is being
asserted.
\end{proof}
